\FloatBarrier
\section{Spectroscopy of Dot Excitations}
\label{app:dot-spectroscopy}
\label{sec:energy-spectroscopy}
\label{sec:measurement-excited-states}

The two-level restriction in Section~\ref{sec:platform-encoding} depends
on the gaps to excluded orbital and valley states. This appendix
explains how those gaps can be inferred from the opening of additional
transport or loading channels, including the effect of spin--valley
mixing near an avoided crossing.

Let $\lvert N,a\rangle$ denote an $N$-electron state, with $a=0$
for the ground state and $a=1$ for the first excited state.
A narrow source--drain window allows the transport cycle
\[
 \lvert N-1,0\rangle
 \xrightarrow{\text{source}\to\text{dot}}\lvert N,0\rangle
 \xrightarrow{\text{dot}\to\text{drain}}\lvert N-1,0\rangle.
\]
Increasing the source--drain bias widens the window and can also
allow the cycle
\[
 \lvert N-1,0\rangle
 \xrightarrow{\text{source}\to\text{dot}}\lvert N,1\rangle
 \xrightarrow{\text{dot}\to\text{drain}}\lvert N-1,0\rangle.
\]
The additional entry process can change the current.

Write $E_{N,a}$ for the energy of $\lvert N,a\rangle$. Starting from an
$(N-1)$-electron state $b$, the energy needed to add an electron
and leave the dot in state $a$ is
\begin{equation}
 \mu_N^{ab}=E_{N,a}-E_{N-1,b}.
 \label{eq:measurement-excited-addition}
\end{equation}
For both entry processes starting from $b=0$, subtraction cancels
the initial-state energy:
\begin{equation}
 \mu_N^{a0}-\mu_N^{00}=E_{N,a}-E_{N,0}\equiv\Delta_a.
 \label{eq:measurement-excitation-gap}
\end{equation}
All energies in this equality refer to the same electrode settings.
This is why we compare transitions with a common initial state
and the same final electron number.

\subsection{Excitation energies from current}
\label{sec:measurement-excitation-current}
At zero temperature and positive bias, use the sequential-tunneling
model with sharp dot levels from
Section~\ref{sec:measurement-sequential-tunneling}.
The narrow-window situation described above satisfies
\[
 \mu_{\mathrm D}<\mu_N^{00}<\mu_{\mathrm S}<\mu_N^{10},
\]
where $a=1$ labels the first excited state. The source has occupied
electron states at the energy $\mu_N^{00}$ needed for ground-state
entry, but none at $\mu_N^{10}$. Hold the drain voltage fixed and
increase $\mu_{\mathrm S}$ by increasing positive $V_b$.
Entry into state $a$ becomes energetically allowed when
$\mu_N^{a0}<\mu_{\mathrm S}$; the equality
$\mu_N^{a0}=\mu_{\mathrm S}$ defines its source-loading boundary.
These addition energies are evaluated at the applied voltages,
since changing the bias can also shift the dot energies.
Figure~\ref{fig:transport-excitation-window}
shows this opening of a second entry process. An electron entering
the excited state may leave directly, or first relax to the ground
state and then leave. Either route can contribute to the current.

\begin{figure}[htbp]
 \centering
 % Schematic transitions from a common (N-1)-electron ground state.
% The horizontal index labels final N-electron eigenstates, not charge.
% The wider window admits the ground state a=0 and excited state a=1.
\begingroup
\small
\setlength{\unitlength}{1mm}
\begin{picture}(150,86)
 \put(15,12){\color{black!9}\rule{119mm}{8mm}}
 \put(15,20){\color{blue!4}\rule{119mm}{32mm}}
 \put(15,52){\color{black!9}\rule{119mm}{26mm}}

 \linethickness{0.8pt}
 \put(15,20){\color{blue!65!black}\line(1,0){119}}
 \put(15,52){\color{blue!65!black}\line(1,0){119}}
 \put(136,17){\makebox(10,6)[l]{$\mu_{\mathrm D}$}}
 \put(136,49){\makebox(10,6)[l]{$\mu_{\mathrm S}$}}
 \put(19,35){\makebox(30,6){bias window}}

 \linethickness{0.5pt}
 \put(15,12){\vector(1,0){124}}
 \put(15,12){\vector(0,1){69}}
 \put(0,78){\makebox(14,6){$\mu_N^{a0}$}}
 \put(141,9){\makebox(7,6){$a$}}
 \multiput(40,12)(26,0){4}{\line(0,-1){1.4}}
 \put(31,4){\makebox(18,6){$0$}}
 \put(57,4){\makebox(18,6){$1$}}
 \put(83,4){\makebox(18,6){$2$}}
 \put(109,4){\makebox(18,6){$3$}}

 % The energy difference compares two final states with the same N.
 \linethickness{0.4pt}
 \color{blue!55!black}
 \multiput(42,30)(2,0){27}{\line(1,0){0.9}}
 \multiput(68,45)(2,0){14}{\line(1,0){0.9}}
 \put(96,30){\line(0,1){15}}
 \put(94.5,30){\line(1,0){3}}
 \put(94.5,45){\line(1,0){3}}
 \color{black}
 \put(99,34.5){\makebox(24,6)[l]{$\Delta_1$}}

 \color{black!42}
 \put(92,62){\circle*{2.4}}
 \put(118,73){\circle*{2.4}}
 \color{blue!65!black}
 \put(40,30){\circle*{2.7}}
 \put(66,45){\circle*{2.7}}
 \color{black}
\end{picture}
\endgroup

 \caption{Bias window for transitions from the $(N-1)$-electron
 ground state to $N$-electron states labeled by $a$. Each point is
 an addition energy $\mu_N^{a0}=E_{N,a}-E_{N-1,0}$.
 The wider window includes the addition energies for the ground
 state ($a=0$) and first excited state ($a=1$). Their difference gives
 $\Delta_1=\mu_N^{10}-\mu_N^{00}=E_{N,1}-E_{N,0}$.
 Gray points mark higher excited-state transitions outside the
 window. Transitions to other charge sectors are excluded.}
 \label{fig:transport-excitation-window}
\end{figure}

To see explicitly why the current changes, consider fast relaxation
within the $N$-electron sector, followed by escape from its ground
state at rate $\Gamma_{\mathrm D}$. Let every escape return the dot
to its $(N-1)$-electron ground state. While the ground-state
transition remains inside the window, the total source entry rate is
\[
 \Gamma_{\mathrm{in}}=
 \begin{cases}
  \Gamma_{\mathrm S0},&\mu_N^{00}<\mu_{\mathrm S}<\mu_N^{10},\\
  \Gamma_{\mathrm S0}+\Gamma_{\mathrm S1},
       &\mu_N^{10}<\mu_{\mathrm S}<\mu_N^{20}.
 \end{cases}
\]
Here $\Gamma_{\mathrm Sa}$ is the entry rate into state $a$ when
that process is allowed. Rates into orthogonal final states add,
as explained in Appendix~\ref{sec:amplitudes-to-master-equation}.
With negligible time spent in the excited
state before relaxation, the rate equation of
Section~\ref{sec:measurement-current-derivation} becomes
\begin{equation}
 \dot p=\Gamma_{\mathrm{in}}(1-p)-\Gamma_{\mathrm D}p,
 \qquad
 I=e\frac{\Gamma_{\mathrm{in}}\Gamma_{\mathrm D}}
             {\Gamma_{\mathrm{in}}+\Gamma_{\mathrm D}},
 \label{eq:measurement-excited-channel-current}
\end{equation}
where $p$ is the probability of $N$ electrons. The new entry rate
therefore produces a current step at
$\mu_N^{10}=\mu_{\mathrm S}$. Finite temperature rounds the step.
If relaxation is slow, the excited-state occupation needs its own
rate equation; the current change can then have a different size
or sign. A state with negligible entry matrix element may produce
no resolved feature.

In the experiment, first locate the ground-state diamond edge and
neighboring charge transitions. Then increase positive $V_b$ in
steps, keeping the drain voltage fixed, and scan $V_g$ at each bias.
Store the current $I(V_g,V_b)$. Taking its numerical bias derivative
$\partial I/\partial V_b$ emphasizes the current steps. Their
positions trace curves in this \emph{differential-conductance map}:
one follows $\mu_N^{00}=\mu_{\mathrm S}$ and another follows
$\mu_N^{a0}=\mu_{\mathrm S}$. We compare this pair of source-alignment
curves, rather than the brightness of arbitrary features in the map.

At a fixed bias, let these two curves occur at $V_g^0$ and $V_g^a$.
If the addition energies shift with the same approximately constant
response $\partial\mu/\partial V_g=-e\alpha$ across their separation,
then
\begin{equation}
 \begin{aligned}
 0&=\mu_N^{a0}(V_g^a)-\mu_N^{00}(V_g^0)
       \simeq\Delta_a-e\alpha(V_g^a-V_g^0),\\
 \Delta_a&\simeq e\alpha(V_g^a-V_g^0).
 \end{aligned}
 \label{eq:measurement-excitation-voltage-spacing}
\end{equation}
Here $\Delta_a$ is evaluated at $V_g^0$, and $\alpha>0$ is
determined from the local voltage response as in
Section~\ref{sec:measurement-bias-spectroscopy}.

Increasing the bias can also permit another electron to enter,
giving $N\to N+1$ and eventually $N+2$. These are different
transitions, not excitations at fixed $N$. For the interpretation
above, other charge sectors must remain inaccessible. In particular,
to forbid further loading from every occupied $N$-electron state
$a$, require
\begin{equation}
 E_{N+1,0}-E_{N,a}
   =\bigl(E_{N+1,0}-E_{N,0}\bigr)-\Delta_a
   >\mu_{\mathrm S}.
 \label{eq:measurement-exclude-further-loading}
\end{equation}
An excited state reduces the cost of adding the next electron, so
checking only the ground-state cost is insufficient. At finite
temperature the margin should be large compared with $k_BT_e$.
Use the neighboring charge boundaries to delimit the scan and
check this condition as excitations become accessible. If other
charge states or initial excited states participate, include their
transitions explicitly before applying
Eq.~\eqref{eq:measurement-excitation-gap}.

\subsection{Excitation energies from loading probabilities}
\label{sec:measurement-excitation-loading}
For a dot coupled to one reservoir, we can instead prepare $N-1$
electrons and ask whether an electron enters during a fixed time.
As the gate lowers successive entry thresholds below the reservoir
chemical potential, more final states become available and loading
becomes faster. No steady current through the target is needed.

Prepare the $(N-1)$-electron ground state before every trial. Let
$V_L$ be the gate voltage during loading and $\Gamma_a$ the bare
entry rate into state $a$. The occupied reservoir states give
\begin{equation}
 \Gamma_{\mathrm{load}}(V_L)
    =\sum_a\Gamma_a f_{\mathrm R}[\mu_N^{a0}(V_L)].
 \label{eq:measurement-loading-rate-spectrum}
\end{equation}
At zero temperature this rate is zero before the first threshold,
$\Gamma_0$ after the ground-state threshold, and
$\Gamma_0+\Gamma_1$ after the first excited-state threshold.
Assume unloading during the loading interval is negligible, either
because the occupied states lie sufficiently below $\mu_{\mathrm R}$
or because rapid internal relaxation traps the loaded electron.
Also exclude further loading using
Eq.~\eqref{eq:measurement-exclude-further-loading} with
$\mu_{\mathrm S}$ replaced by $\mu_{\mathrm R}$.

If $q(t)$ is the probability that no electron has yet entered, then
$\dot q=-\Gamma_{\mathrm{load}}q$ with $q(0)=1$. The probability
of finding $N$ electrons after a loading time $\tau_L$ is therefore
\begin{equation}
 P_N(\tau_L)=1-e^{-\Gamma_{\mathrm{load}}\tau_L},
 \qquad
 \Gamma_{\mathrm{load}}=-\frac{\ln[1-P_N(\tau_L)]}{\tau_L}.
 \label{eq:measurement-pulsed-spectroscopy}
\end{equation}
To measure $P_N$, pulse the gate to $V_L$, wait $\tau_L$, then
close the reservoir barrier to retain the charge during sensor
readout. The fraction of repeated trials reporting $N$ electrons
estimates $P_N$. Repeat at different $V_L$ and loading times;
choose times for which the probability is not saturated at one.
Equation~\eqref{eq:measurement-pulsed-spectroscopy} then converts
the counts into a loading-rate curve. Its successive onsets give
the threshold voltages~\cite{Thalakulam2011}.

If unloading is appreciable, extract the rate from the initial
growth instead. For times short compared with the tunneling
waiting times, the rate description gives
\begin{equation}
 P_N(V_L,\tau_L)\simeq\Gamma_{\mathrm{load}}(V_L)\tau_L.
 \label{eq:measurement-short-time-loading}
\end{equation}
Entry followed by exit requires two events and contributes only
at higher order in $\tau_L$. Repeated measurements at several
short loading times therefore give $\Gamma_{\mathrm{load}}$
from the initial slope.

Two onsets aligned with the same $\mu_{\mathrm R}$ give a gap by
the same subtraction as
Eq.~\eqref{eq:measurement-excitation-voltage-spacing}, using the
voltage response for this loading configuration. Without an
independent voltage-to-energy conversion, this experiment determines
threshold separations in volts, not absolute gaps in energy units.

\subsection{Spin--valley mixing}
\label{sec:measurement-spin-valley}

Consider loading one electron into a dot whose orbital excitations
lie outside the energy range being probed. The valley eigenstates
$\ket{v_-}$ and $\ket{v_+}$ include the interface-induced valley coupling
described in Section~\ref{sec:device-orbital-valley}. Their
separation is $\Delta_v>0$. At zero field, each valley has two
degenerate spin states. As the magnetic field $B$ increases, $\ket{v_-,\uparrow}$
rises toward $\ket{v_+,\downarrow}$. A spin-dependent coupling
mixes this pair and prevents their energies from crossing.
We measure the minimum separation to determine that coupling.

Take the field along $z$, with $B\geq0$ and the same
$g>0$ in both valleys. Neglect changes of the orbital and valley
energies with field. Before including spin--valley mixing, the
four energies, measured relative to $\ket{v_-,\downarrow}$,
are\footnote{This four-state description is specific to one
electron. For several electrons, exchange can separate different
total-spin multiplets already at $B=0$, and one energy can contain
more than one multiplet.}
\begin{equation}
 \begin{array}{c|cc}
  &\downarrow&\uparrow\\ \hline
  v_-&0&E_Z\\
  v_+&\Delta_v&\Delta_v+E_Z
 \end{array},
 \qquad E_Z=g\mu_BB.
 \label{eq:measurement-spin-valley-levels}
\end{equation}
The zero-field valley gap is obtained by applying the loading
spectroscopy of Section~\ref{sec:measurement-excitation-loading} to
an empty dot. If $\mu_-$ and $\mu_+$ are the addition energies
into the two valleys, their difference is
\begin{equation}
 \Delta_v=\mu_+-\mu_-.
 \label{eq:measurement-valley-addition-gap}
\end{equation}
The two degenerate spins in each valley contribute to the loading
rate at the same energy. Under the assumed exclusion of orbital
excitations, these two distinct loading energies correspond to
the two valleys.

\subsection{The spin--valley avoided crossing}
As $B$ increases, the spin excitation energy $E_Z$ approaches
the valley gap $\Delta_v$, bringing $\ket{v_-,\uparrow}$ and
$\ket{v_+,\downarrow}$ close in energy. Without coupling, the two energies
would cross at $g\mu_BB=\Delta_v$. A spin-dependent interaction,
such as spin--orbit interaction, can couple these states.
Denote this interaction by $V_{\mathrm{sv}}$ and its matrix
element by
\[
 \lambda=\bra{v_-,\uparrow}V_{\mathrm{sv}}\ket{v_+,\downarrow}.
\]

To describe the crossing, retain this pair of nearly degenerate
states. The two remaining states lie about $\Delta_v$ below
and above the pair. Couplings to them therefore give small
perturbative corrections when their matrix elements are much
smaller than $\Delta_v$. The other opposite-spin,
opposite-valley pair, $\ket{v_-,\downarrow}$ and
$\ket{v_+,\uparrow}$, can also couple, but its separation is
$\Delta_v+E_Z$, rather than the small difference
$\Delta_v-E_Z$. Neglecting the off-resonant corrections and
diagonal shifts from $V_{\mathrm{sv}}$, we obtain, in the basis
$(\ket{v_-,\uparrow},\ket{v_+,\downarrow})$,
\begin{equation}
 H_{\mathrm{cross}}=
 \begin{pmatrix}E_Z&\lambda\\\lambda^*&\Delta_v\end{pmatrix},
 \qquad
 \Delta_\pm=\frac{E_Z+\Delta_v}{2}
       \pm\sqrt{\frac{(E_Z-\Delta_v)^2}{4}+|\lambda|^2}.
 \label{eq:measurement-spin-valley-anticrossing}
\end{equation}
Here $\Delta_-$ and $\Delta_+$ are the excitation energies
of the two mixed states, relative to the unperturbed
$\ket{v_-,\downarrow}$ ground state~\cite{Yang2013SpinValley}.
Their separation is
\begin{equation}
 d(B)=\Delta_+-\Delta_-
     =\sqrt{(g\mu_BB-\Delta_v)^2+4|\lambda|^2}.
 \label{eq:measurement-spin-valley-separation}
\end{equation}
For nonzero $\lambda$ the levels remain separated even at the
field where the uncoupled energies coincide: this is the
\emph{avoided crossing}. If $g$, $\Delta_v$, and $\lambda$
are approximately constant over the field range,
\begin{equation}
 B_*=\frac{\Delta_v}{g\mu_B},
 \qquad d(B_*)=\min_B d(B)=2|\lambda|.
 \label{eq:measurement-spin-valley-minimum-gap}
\end{equation}
Measuring this minimum separation therefore determines
the magnitude of the coupling.

\subsection{Extracting the coupling from loading spectra}
In loading spectroscopy, an electron can enter the empty dot
and occupy either mixed state. The two processes have the same
initial state, so their addition-energy difference equals
$\Delta_+-\Delta_-$, as in
Section~\ref{sec:measurement-excitation-loading}. We use this
difference to measure $d(B)$.

At each fixed $B$, use that section's $N=1$ loading protocol,
keeping entry of a second electron forbidden. Scan the loading
gate voltage $V_L$ and obtain the loading-rate curve from
Eq.~\eqref{eq:measurement-short-time-loading}.
After the ground-state loading step, entry into each mixed
state adds a further step. Fit these steps with the Fermi
functions in Eq.~\eqref{eq:measurement-loading-rate-spectrum}
to locate the two alignment voltages: $V_{L,1}(B)$ for the
state at $\Delta_-$ and $V_{L,2}(B)$ for the state at $\Delta_+$.

Suppose both addition energies have the same approximately
constant response $\partial\mu/\partial V_L=-e\alpha_L$,
with $\alpha_L>0$, across this interval. Then
Eq.~\eqref{eq:measurement-excitation-voltage-spacing} gives
\begin{equation}
 d(B)\simeq e\alpha_L\bigl[V_{L,2}(B)-V_{L,1}(B)\bigr].
 \label{eq:measurement-spin-valley-voltage-gap}
\end{equation}
If $\alpha_L$ is not already known, an independently measured
spin gap can provide an energy reference. At a field where
$E_Z\ll\Delta_v$ is resolved, the spin-resonance measurement of
Section~\ref{sec:measurement-rabi} gives a frequency
$f_{\mathrm{res}}$ with $hf_{\mathrm{res}}=E_Z$.
Use the loading measurement at that field to find the voltage
separation $\delta V_L^s$ between $\ket{v_-,\downarrow}$ and
$\ket{v_-,\uparrow}$. Their known energy separation gives
\begin{equation}
 e\alpha_L\simeq\frac{E_Z}{\delta V_L^s}
              =\frac{hf_{\mathrm{res}}}{\delta V_L^s}.
 \label{eq:measurement-spin-loading-calibration}
\end{equation}
Using this response for the crossing measurement requires it
to remain approximately unchanged over the relevant gate and
field range. The common-response approximation also requires
the gap to vary negligibly between the two loading voltages.

Repeat the loading scan at fields on both sides of $B_*$,
with the same gate settings apart from the scan. Convert each
voltage separation to $d(B)$ and compare the resulting curve
with Eq.~\eqref{eq:measurement-spin-valley-separation}.
Its minimum gives $2|\lambda|$; the points on either side
check the predicted field dependence. Both loading steps must
be detectable, and their minimum separation must exceed the
experimental energy resolution for this direct extraction.

The measured coupling quantifies the departure from a pure
lower-valley spin state. Away from the crossing, where
$|\lambda|\ll|\Delta_v-E_Z|$, first-order perturbation theory
gives
\begin{equation}
 \ket{\widetilde{\uparrow}}\simeq
 \ket{v_-,\uparrow}
 -\frac{\lambda^*}{\Delta_v-E_Z}\ket{v_+,\downarrow},
 \qquad
 w_v\simeq\frac{|\lambda|^2}{(\Delta_v-E_Z)^2}.
 \label{eq:measurement-spin-valley-admixture}
\end{equation}
Here $w_v$ is the leading higher-valley weight due to this
coupling, not a gate-error probability. It estimates how far
from the crossing the qubit should operate to keep that
admixture small, using $\lambda$ at the operating field or
assuming that its variation from the measured crossing is
negligible.
